\section{Related work}\label{sec:related}

\subsection{Energy and variational training of neural solvers}
The deep Ritz method \cite{EYu2018} minimises the continuous energy of a variational problem
over neural trial functions, and the deep energy method \cite{Samaniego2020} brought the idea
to computational mechanics; later work treated essential boundary conditions
\cite{LiaoMing2021} and the error of the quadrature that evaluates the continuous functional
\cite{LiuCaiRamani2023}, and variational and weak-form losses followed
\cite{Kharazmi2021,Wang2022CENN,Wang2025KINN,Manav2024}. Physics-informed neural networks
\cite{Raissi2019} minimise instead a Euclidean residual of the differential equation. Robust
variational physics-informed networks \cite{Rojas2024} and the deep Fourier residual method
\cite{Taylor2023DFR,Taylor2024} construct losses equivalent to the energy norm of the error,
whose discrete form is the energy norm of \cref{sec:theory}, by measuring the residual in a
dual norm; \cref{lem:exact} and \cref{rem:norms} are the exact algebraic counterpart of that
equivalence on the assembled system, without quadrature error or equivalence constants,
because the inverse in the dual norm cancels. Energy natural gradient descent trains
physics-informed networks in the geometry that the energy induces on the parameters
\cite{MullerZeinhofer2023}.

\subsection{Training without solutions, and its comparison with supervision}
Surrogates trained on a physical objective instead of solutions have a history in this
journal: on a variational objective for elliptic problems with random inputs \cite{Zhu2019},
on the residual of autoregressive dynamics \cite{GenevaZabaras2020} and on irregular domains
mapped to a reference grid \cite{Gao2021PhyGeoNet}; operators have been trained the same way
\cite{Li2024PINO,WangWangPerdikaris2021}. Variational operator learning \cite{Xu2024VOL}
trains operators on the finite-element system, evaluated matrix-free, and uses the gradient
identity of \cref{lem:exact} as a computational device; its loss is the Euclidean norm of
the residual, or a regression on provisional labels produced by a few conjugate-gradient
iterations. The variational physics-informed neural operator \cite{Eshaghi2025VINO}
evaluates the continuous functional analytically on its own bilinear discretisation over
uniform grids. Larson et~al.~\cite{Larson2026MLPEnergy} train multilayer perceptrons that
map the data of parametrised, possibly nonlinear, problems to the nodal values of their
finite-element solutions by minimising the energy assembled element by element, without
precomputed solutions. Other label-free approaches build on the finite-element residual or
weak form \cite{Lee2025FEONet,Uriarte2022,Gao2022,Goswami2022}. Closest to the present work
are graph networks on finite-element meshes trained by minimising the potential energy
\cite{Dalton2023}, physics-informed MeshGraphNets trained without data for non-stationary
and nonlinear problems \cite{Wurth2024}, a three-dimensional convolutional network for
elastic fields in periodic polycrystals, trained without reference solutions on the
divergence of its predicted stress \cite{MonteiroFernandes2025}, and the pretrain finite
element method \cite{Wang2026PretrainFEM}, which trains a point-cloud transformer operator
on the potential energy, with derivatives from finite-element shape functions, and uses its
prediction to warm-start the finite-element solver; a follow-up applies that approach to
three-dimensional pressure vessels \cite{Sun2026FEVessel}. Several of these works compare
label-free with data-driven training. Physics-constrained surrogates trained on a
variational objective gave predictions comparable to data-driven ones \cite{Zhu2019}. Three
train the same network on the same data with and without reference solutions, in single
runs. In soft-tissue beams the two had comparable displacement errors and the label-free one
smaller errors in energy and deformation gradient, a gap that adding the deformation
gradient to the supervised loss narrowed but did not close \cite{Dalton2023}; in
polycrystals the label-free network was more accurate in stress than the supervised one,
which was trained on the reference stresses \cite{MonteiroFernandes2025}; and in pressure
vessels the supervised one was more accurate in displacement and in the $H^1$ seminorm
\cite{Sun2026FEVessel} (\cref{sec:discussion}). For finite-element surrogates we are not
aware of a comparison that changes the norm of a supervised loss while keeping its labels,
nor of a comparison with and without labels repeated over several seeds under criteria fixed
in advance.

\subsection{The norm of the loss}
Supervised mesh surrogates, from message-passing networks \cite{Pfaff2021,Dalton2022} to
operator networks \cite{Lu2021DeepONet,Li2021FNO} and geometry-aware operators
\cite{He2024GeomDeepONet,Zeng2025PCNO,Liu2026GINOT}, are typically trained and evaluated
with relative $L^2$ errors of the predicted fields. That the norm of a loss shapes what a
network learns is known in machine learning: training on derivatives as well as values,
Sobolev training, improves the accuracy of the derivatives \cite{Czarnecki2017}, and for
physics-informed networks the $L^2$ residual loss can fail to control the solution where
another norm does \cite{WangL2NeurIPS2022}. For finite-element surrogates of linear
elasticity the right norm is fixed by the problem itself, it is the energy norm, and the
energy supplies it without solutions. \Cref{rem:rank} and \cref{sec:checks:ranking} show
that the displacement error can rank surrogates in the opposite order to their stress
errors.

\subsection{Error estimation without the solution}
A posteriori error estimation bounds the energy-norm error of a finite-element approximation
without the exact solution \cite{AinsworthOden2000}: from a statically admissible stress
field through the hypercircle of Prager and Synge \cite{PragerSynge1947}, or through the
error in the constitutive relation \cite{LadevezeLeguillon1983}, which has been applied to
the solutions of physics-informed networks in elasticity, with the admissible pair of
displacement and stress given by a network of mixed form \cite{GuoHaghighat2022}. Krylov
methods estimate the energy-norm error of an iterate by Gauss quadrature
\cite{GolubMeurant1997,StrakosTichy2002}. The checks of \cref{sec:theory:detector} are
weaker: they rank predictions and detect gross failures, but need nothing beyond the
assembled $K$ and $F$. The nearest uses we found are a residual-norm test that decides
whether a warm-started solve is needed \cite{Wang2026PretrainFEM} and an energetic upper
bound on the effective stiffness computed from predicted fields
\cite{MonteiroFernandes2025}.

\subsection{Solvers and reduced bases}
Neural predictions have been used as initial iterates of iterative solvers
\cite{Wang2026PretrainFEM,Eshaghi2026NOWS,Sun2026FEVessel} (supplementary
Section~\SuppWarm{}), inside hybrid iterative schemes \cite{Lee2025MetaSolver}, and to enrich
finite-element approximation spaces \cite{Barucq2026Enrich}. Reduced basis methods, the
classical surrogates of parametrised linear problems \cite{Hesthaven2016RB}, work in one
truth space, usually on a reference domain to which each geometry is mapped, and rely for
efficiency on an affine, or affinely approximated, dependence on the parameters; the
instances here differ in the number of holes or cavities, which no such mapping covers, and
in their meshes. The Ritz projection of \cref{sec:methods:ritz} is a classical approximation
that does apply to them, at the price of the stiffness matrix of each instance.
